\documentclass[conference]{IEEEtran}
\IEEEoverridecommandlockouts

\usepackage{cite}
\usepackage{amsmath,amssymb,amsfonts}
\usepackage{algorithmic}
\usepackage{graphicx}
\usepackage{textcomp}
\usepackage{xcolor}
\usepackage{hyperref}
\usepackage{booktabs}

\begin{document}

\title{Unsupervised Statistical Learning for Kinematic Anomaly Detection in Autonomous Driving Scenarios}

\author{\IEEEauthorblockN{Pavel Stepanov}
\IEEEauthorblockA{\textit{Dept. of Computer Science} \\
\textit{University of Miami}\\
Coral Gables, USA \\
pas273@miami.edu}
\and
\IEEEauthorblockN{Ramses Loaces}
\IEEEauthorblockA{\textit{Dept. of Computer Science} \\
\textit{University of Miami}\\
Coral Gables, USA \\
rxl1238@miami.edu}
\and
\IEEEauthorblockN{Justin Cabrera}
\IEEEauthorblockA{\textit{Dept. of Computer Science} \\
\textit{University of Miami}\\
Coral Gables, USA \\
justin.cabrera2718@gmail.com}
}

\maketitle

\begin{abstract}
Autonomous vehicles must operate safely in environments characterized by stochastic human behavior. While nominal driving scenarios are well-represented in current datasets, identifying anomalous events---such as aggressive maneuvering or near-miss interactions---remains a challenge due to the scarcity of labeled high-risk data. This paper presents an unsupervised statistical learning framework applied to the Waymo Open Motion Dataset (v1.3.1). We extract kinematic features from over 10,000 traffic agents and utilize Principal Component Analysis (PCA) for dimensionality reduction. We systematically compare parametric (Mahalanobis Distance) and non-parametric (One-Class Support Vector Machine) approaches for outlier detection. Experimental results demonstrate that statistically identified anomalies correlate strongly with ground-truth interaction flags, validating the efficacy of unsupervised methods for filtering safety-critical scenarios.
\end{abstract}

\begin{IEEEkeywords}
Statistical Learning, Anomaly Detection, Waymo Open Dataset, Unsupervised Learning, PCA, One-Class SVM
\end{IEEEkeywords}

\section{Introduction}
The deployment of autonomous vehicles (AVs) requires rigorous validation against the "long tail" of driving distributions. Traditional supervised learning approaches are often limited by the scarcity of labeled anomaly data and the subjective nature of what constitutes an "unsafe" maneuver. Consequently, unsupervised novelty detection techniques are essential for mining large-scale datasets for safety-critical events \cite{pimentel2014review}. This study applies classical statistical learning methods to quantify kinematic deviance in urban traffic scenarios using the high-fidelity Waymo Open Motion Dataset \cite{ettinger2021waymo}.

\section{Data Description}
We utilize the Waymo Open Motion Dataset (WOMD) v1.3.1, a large-scale collection of object trajectories mined from real-world urban environments \cite{sun2020scalability}. The dataset provides 20Hz object tracks sampled over 20-second segments, including rich metadata such as agent type (vehicle, pedestrian, cyclist) and interaction connectivity. Our study focuses on the validation split, processing raw Protocol Buffer (TFRecord) data to extract a structured tabular dataset consisting of $N \approx 10,000$ agents across 300 distinct scenarios.

\subsection{Feature Extraction}
For each agent, we extracted kinematic features at the prediction horizon ($t=0$) to capture the instantaneous state of the agent during peak interaction density. The features include:
\begin{itemize}
    \item \textbf{Speed ($v$):} Magnitude of the velocity vector ($m/s$).
    \item \textbf{Acceleration ($a$):} Derived from the rate of change of speed ($m/s^2$).
    \item \textbf{Yaw Rate ($\dot{\psi}$):} Angular velocity representing the rate of change of heading ($rad/s$).
    \item \textbf{Lateral G-Force:} Computed as $|v \times \dot{\psi}|$, serving as a proxy for cornering intensity.
\end{itemize}

\section{Methodology}

\subsection{Preprocessing}
Data cleaning involved the removal of the autonomous vehicle (SDC) tracks to focus exclusively on human-driven agent behavior. Features were normalized using Z-score standardization to ensure that units of different magnitudes (e.g., speed vs. yaw rate) contributed equally to distance calculations.

\subsection{Dimensionality Reduction}
We applied Principal Component Analysis (PCA) \cite{jolliffe2002principal} to reduce the dimensionality of the feature space. By projecting the four-dimensional kinematic vector onto the first two principal components, we retained over 85\% of the total variance, allowing for the visualization of "nominal" clusters and isolated outliers.

\subsection{Anomaly Detection}
Two unsupervised methods were implemented to identify deviations from standard driving profiles:
\begin{enumerate}
    \item \textbf{Mahalanobis Distance:} A multivariate distance metric that accounts for the covariance structure of the data \cite{mahalanobis1936generalized}. Observations exceeding the $\chi^2$ critical value at a 99\% confidence interval were flagged as statistical outliers.
    \item \textbf{One-Class SVM (OC-SVM):} A non-parametric kernel method designed to estimate the support of a high-dimensional distribution \cite{scholkopf2001estimating}. We utilized a Radial Basis Function (RBF) kernel with a nu-parameter ($\nu=0.05$) to define a robust boundary around the majority of data points.
\end{enumerate}

\section{Results}
\subsection{Statistical Validation}
A key feature of the Waymo dataset is the \texttt{IsInteractive} flag, which identifies agents involved in complex maneuvers with other road users \cite{ettinger2021waymo}. We validated our detected anomalies against this flag. A contingency table analysis reveals that agents exhibiting extreme Mahalanobis distances are significantly more likely to be involved in interactions compared to the baseline population. This suggests that "anomalous" kinematics in a statistical sense are often proxies for "safety-critical" events in a physical sense.

\section{Conclusion}
Our analysis confirms that unsupervised statistical learning is a viable methodology for identifying high-risk driving behaviors without the need for manual annotation. The integration of PCA for feature compression and Mahalanobis Distance for outlier scoring provides a computationally efficient pipeline for filtering the "long tail" of the Waymo Open Motion Dataset. Future work will explore the temporal evolution of these anomalies across the full 20-second trajectory.

\bibliographystyle{IEEEtran}
\bibliography{references}

\end{document}